\usepackage{tcolorbox}
\tcbuselibrary{most}

\newtcbtheorem{examct}%
{考试内容}{colback=Emerald!10,colframe=cyan!40!black,fonttitle=\bfseries}{exact}
\newtcbtheorem{examrqm}%
{考试要求}{colback=Emerald!10,colframe=cyan!40!black,fonttitle=\bfseries}{exarqm}

\newtcbtheorem[number within=section]{defn}%
{定义}{colback=OliveGreen!10,colframe=Green!70,fonttitle=\bfseries}{def}

\newtcbtheorem[number within=section]{lemma}%
{引理}{colback=Salmon!20,colframe=Salmon!90!Black,fonttitle=\bfseries}{lem}

\newtcbtheorem[number within=section]{prove}%
{证明}{colback=TealBlue!20,colframe=TealBlue!90!Black,fonttitle=\bfseries}{pro}

\newtcbtheorem{remark}{注意}{
  colback=gray!10,   % 浅灰背景
  colframe=white,    % 白色边框（视觉上无边框）
  fonttitle=\bfseries,
  boxrule=0pt,       % 边框粗细为0（彻底无边框）
  arc=0pt,           % 直角
  before upper={\parindent=2em}  % 内容首行缩进2字符
}{rem}
% 使用另一个计数器 use counter from=lemma
\newtcbtheorem[use counter from=lemma, number within=section]{them}%
{定理}{colback=SeaGreen!10!CornflowerBlue!10,colframe=RoyalPurple!55!Aquamarine!100!,fonttitle=\bfseries}{them}

\newtcbtheorem[use counter from=lemma, number within=section]{criterion}%
{准则}{colback=green!5,colframe=green!35!black,fonttitle=\bfseries}{cri}

\newtcbtheorem[use counter from=lemma, number within=section]{corollary}%
{推论}{colback=Emerald!10,colframe=cyan!40!black,fonttitle=\bfseries}{cor}


% colback=red!5,colframe=red!75!black
